% Figure 2: Complex-valued U-Net architecture
% Compile standalone or \input{} into main document
% Requires: \usepackage{tikz}
%           \usetikzlibrary{arrows.meta, positioning, calc, fit, backgrounds, decorations.pathreplacing}

\begin{figure*}[t]
\centering
\begin{tikzpicture}[
    >=Stealth,
    node distance=0.6cm and 0.8cm,
    % Styles
    enc/.style={draw, fill=blue!15, rounded corners=2pt, minimum width=1.1cm, align=center,
                font=\tiny},
    dec/.style={draw, fill=orange!12, rounded corners=2pt, minimum width=1.1cm, align=center,
                font=\tiny},
    mid/.style={draw, fill=violet!12, rounded corners=2pt, minimum width=1.3cm, align=center,
                font=\tiny},
    io/.style={draw, fill=green!10, rounded corners=2pt, minimum width=1.4cm, align=center,
               font=\tiny},
    aux/.style={draw, fill=yellow!15, rounded corners=2pt, minimum width=1.2cm, align=center,
                font=\tiny},
    skip/.style={->, gray!60, thick, dashed},
    arrow/.style={->, thick},
    label/.style={font=\tiny, text=gray!70!black},
    darrow/.style={->, thick, dashed},
]

% ===== INPUTS (left) =====
\node[io] (er) at (-0.5, 0) {$E_r, E_i$\\(complex field)};
\node[aux] (helm) at (-0.5, -1.3) {Helmholtz\\residual $H_r, H_i$};
\node[aux] (auxin) at (-0.5, -2.6) {$\varepsilon$, $m_\text{src}$\\(real aux)};

% RealToComplex bridge
\node[aux, fill=yellow!25] (r2c) at (1.8, -2.6) {RealTo-\\Complex};
\draw[arrow] (auxin) -- (r2c);

% Stem
\node[enc, minimum height=1.2cm, minimum width=1.3cm, fill=blue!8] (stem) at (3.2, -1.0)
    {Stem\\ComplexConv\\$\to$ 40\,ch};
\draw[arrow] (er) -| (stem);
\draw[arrow] (helm) -| (stem);
\draw[arrow] (r2c) -| (stem);

% ===== ENCODER =====
% Level 1 (full res, ds=1): 40 complex channels
\node[enc, minimum height=1.8cm] (e1) at (5.0, -1.0)
    {CResBlock\\$\times 3$\\[2pt]40\,ch};
\draw[arrow] (stem) -- (e1);

% Downsample 1
\node[enc, fill=blue!25, minimum height=0.5cm, font=\tiny] (d1) at (6.2, -1.0) {$\downarrow 2$};
\draw[arrow] (e1) -- (d1);

% Level 2 (ds=2): 80 complex channels
\node[enc, minimum height=1.8cm] (e2) at (7.4, -1.0)
    {CResBlock\\$\times 3$\\[2pt]80\,ch};
\draw[arrow] (d1) -- (e2);

% Downsample 2
\node[enc, fill=blue!25, minimum height=0.5cm, font=\tiny] (d2) at (8.6, -1.0) {$\downarrow 2$};
\draw[arrow] (e2) -- (d2);

% Level 3 (ds=4): 160 complex channels
\node[enc, minimum height=1.8cm] (e3) at (9.8, -1.0)
    {CResBlock\\$\times 3$\\[2pt]160\,ch};
\draw[arrow] (d2) -- (e3);

% Downsample 3
\node[enc, fill=blue!25, minimum height=0.5cm, font=\tiny] (d3) at (11.0, -1.0) {$\downarrow 2$};
\draw[arrow] (e3) -- (d3);

% Level 4 (ds=8): 320 complex channels + 8-head attention
\node[enc, minimum height=2.0cm] (e4) at (12.2, -1.0)
    {CResBlock\\$\times 3$\\+ CAttn (8h)\\320\,ch};
\draw[arrow] (d3) -- (e4);

% ===== BOTTLENECK =====
\node[mid, minimum height=2.2cm, minimum width=1.5cm] (bot) at (12.2, -4.0)
    {CResBlock\\CAttn (8h)\\CResBlock\\320\,ch};
\draw[arrow] (e4) -- (bot);

% ===== DECODER =====
% Level 4 (decoder): 320 ch + attention
\node[dec, minimum height=2.0cm] (de4) at (12.2, -7.0)
    {CResBlock\\$\times 4$\\+ CAttn (8h)\\320\,ch};
\draw[arrow] (bot) -- (de4);

% Upsample 3
\node[dec, fill=orange!25, minimum height=0.5cm, font=\tiny] (u3) at (11.0, -7.0) {$\uparrow 2$};
\draw[arrow] (de4) -- (u3);

% Level 3 (decoder): 160 ch
\node[dec, minimum height=1.8cm] (de3) at (9.8, -7.0)
    {CResBlock\\$\times 4$\\160\,ch};
\draw[arrow] (u3) -- (de3);

% Upsample 2
\node[dec, fill=orange!25, minimum height=0.5cm, font=\tiny] (u2) at (8.6, -7.0) {$\uparrow 2$};
\draw[arrow] (de3) -- (u2);

% Level 2 (decoder): 80 ch
\node[dec, minimum height=1.8cm] (de2) at (7.4, -7.0)
    {CResBlock\\$\times 4$\\80\,ch};
\draw[arrow] (u2) -- (de2);

% Upsample 1
\node[dec, fill=orange!25, minimum height=0.5cm, font=\tiny] (u1) at (6.2, -7.0) {$\uparrow 2$};
\draw[arrow] (de2) -- (u1);

% Level 1 (decoder): 40 ch
\node[dec, minimum height=1.8cm] (de1) at (5.0, -7.0)
    {CResBlock\\$\times 4$\\40\,ch};
\draw[arrow] (u1) -- (de1);

% ===== OUTPUT =====
\node[io] (out) at (3.2, -7.0) {CGNorm\\ModReLU\\CConv $\to 1$};
\draw[arrow] (de1) -- (out);

\node[io, fill=orange!20] (vout) at (1.0, -7.0) {$v_r, v_i$\\(velocity)};
\draw[arrow] (out) -- (vout);

% ===== SKIP CONNECTIONS =====
\draw[skip] (e1.south) -- ++(0, -1.0) -| (de1.north);
\draw[skip] (e2.south) -- ++(0, -0.8) -| (de2.north);
\draw[skip] (e3.south) -- ++(0, -0.6) -| (de3.north);
\draw[skip] (e4.south west) -- ++(-0.2, -0.5) |- (de4.north west);

% Skip labels
\node[label] at (5.0, -3.0) {skip};
\node[label] at (7.4, -3.2) {skip};
\node[label] at (9.8, -3.4) {skip};

% ===== TIME/COND EMBEDDING =====
\node[aux, fill=red!10, minimum width=2.0cm] (temb) at (8.0, -4.5)
    {$t$ sinusoidal emb\\$+\;\lambda_\text{norm}$\\$\to$ MLP $\to$ emb};

% Arrows from temb to all resblocks (simplified: one arrow to encoder, one to decoder)
\draw[darrow, red!40, ->] (temb.west) -- ++(-2.0, 0) node[label, above, pos=0.5] {modulates all CResBlocks};
\draw[darrow, red!40, ->] (temb.east) -- ++(2.0, 0);

% ===== INSET: CResBlock =====
\begin{scope}[shift={(-0.5, -4.8)}, scale=0.85]
    \node[font=\tiny\bfseries] at (0.8, 0.3) {ComplexResBlock:};
    \node[enc, minimum width=1.8cm, minimum height=0.35cm] (r1) at (0.8, -0.1) {CGNorm $\to$ CConv};
    \node[enc, minimum width=1.8cm, minimum height=0.35cm, fill=green!15] (r2) at (0.8, -0.6) {ModReLU};
    \node[enc, minimum width=1.8cm, minimum height=0.35cm, fill=red!8] (r3) at (0.8, -1.1)
        {$(s_r{+}is_i) \cdot h + h$};
    \node[enc, minimum width=1.8cm, minimum height=0.35cm] (r4) at (0.8, -1.6) {CGNorm $\to$ ModReLU};
    \node[enc, minimum width=1.8cm, minimum height=0.35cm] (r5) at (0.8, -2.1) {CConv (zero-init)};
    \draw[arrow, thin] (r1) -- (r2);
    \draw[arrow, thin] (r2) -- (r3);
    \draw[arrow, thin] (r3) -- (r4);
    \draw[arrow, thin] (r4) -- (r5);
    % skip
    \draw[skip, thin] (-0.5, -0.1) -- (-0.5, -2.1) -- (r5.west);
    \node[label] at (-0.7, -1.1) {+};
    % emb annotation
    \node[font=\tiny, text=red!60!black] at (2.2, -1.1) {$\leftarrow$ emb};
\end{scope}

\end{tikzpicture}
\caption{Complex-valued U-Net architecture. The encoder (blue) downsamples through four resolution levels with complex
residual blocks (CResBlock) and complex self-attention (CAttn) at the coarsest level. Real auxiliary inputs
($\varepsilon$, source mask) are projected into complex space via a RealToComplex bridge. The Helmholtz residual of the
current field is concatenated as an additional complex input channel. The timestep $t$ and normalized wavelength
$\lambda$ are embedded and injected into every residual block via complex scale modulation. Skip connections (dashed
gray) feed encoder features to the decoder (orange). The inset (bottom left) shows the internal structure of each
ComplexResBlock: complex group normalization (CGNorm), complex convolution (CConv), ModReLU activation, embedding
modulation, and a zero-initialized output convolution with residual addition.}
\label{fig:architecture}
\end{figure*}
